\documentclass[11pt]{article}

\usepackage[utf8]{inputenc}
\usepackage[T1]{fontenc}
\usepackage[spanish]{babel}
\usepackage[a4paper,landscape,left=0cm,right=0cm,top=0.2cm,bottom=0cm]{geometry}
\usepackage{tikz}
\usepackage{xcolor}
\usepackage{graphicx}
\usepackage{lmodern}
\pagestyle{empty}
\setlength{\parindent}{0pt}

\definecolor{azulUVE}{RGB}{0,120,180}
\definecolor{naranjaUVE}{RGB}{220,90,20}

\begin{document}
\thispagestyle{empty}
\centering

\resizebox{\paperwidth}{!}{%
\begin{tikzpicture}[x=1cm,y=1cm]

% Caja total del dibujo
\path[use as bounding box] (0,0) rectangle (27,18);

%------------------ TITULO ------------------
\node[align=center] at (13.5,17.1) {
    {\LARGE\bfseries\textcolor{azulUVE}{V de Gowin para estadística}}\\[0.18cm]
    {\normalsize\bfseries Tema: Análisis de la relación entre el uso de redes sociales y el rendimiento académico en estudiantes universitarios}
};

%------------------ ENCABEZADOS ------------------
\node[anchor=west, align=left] at (0.7,15.6) {
    {\large\bfseries\textcolor{azulUVE}{PARTE CONCEPTUAL}}\\[0.05cm]
    \textbf{Pienso sobre el tema}
};

\node[anchor=west, align=left] at (19.6,15.6) {
    {\large\bfseries\textcolor{azulUVE}{PARTE METODOLÓGICA}}\\[0.05cm]
    \textbf{Actúo sobre el tema}
};

%------------------ LINEAS DE LA V ------------------
\draw[line width=1pt] (7.6,14.6) -- (13.5,3.0);
\draw[line width=1pt] (19.4,14.6) -- (13.5,3.0);

%------------------ 1 PREGUNTA ------------------
\node[
    text width=6.9cm,
    align=center,
    font=\small
] at (13.5,13.8) {
    {\fontsize{22}{22}\selectfont\bfseries\textcolor{azulUVE}{1}}\\
    {\large\bfseries Pregunta clave}\\[0.04cm]
    {\color{naranjaUVE}\bfseries ¿Qué queremos conocer?}\\[0.16cm]
    ¿Existe una relación entre el uso de redes sociales y el rendimiento académico en estudiantes universitarios, considerando la distracción académica, las horas de estudio y la asistencia?
};

%------------------ 2 OBJETO ------------------
\node[
    text width=8.9cm,
    align=center,
    font=\small
] at (13.5,1.95) {
    {\fontsize{22}{22}\selectfont\bfseries\textcolor{azulUVE}{2}}\\
    {\large\bfseries Objeto de estudio}\\[0.04cm]
    {\color{naranjaUVE}\bfseries ¿Qué fenómeno estudiamos?}\\[0.16cm]
    La asociación entre el tiempo diario en redes sociales y el promedio académico de estudiantes universitarios, tomando en cuenta factores académicos como distracción, horas de estudio y asistencia.
};

%------------------ LADO IZQUIERDO ------------------

% 3 Conceptos
\node[
    anchor=north west,
    text width=7.25cm,
    align=left,
    font=\fontsize{8.4}{10.1}\selectfont
] at (0.7,15) {
    {\fontsize{22}{22}\selectfont\bfseries\textcolor{azulUVE}{3}} {\large\bfseries Conceptos relacionados}\\
    {\color{naranjaUVE}\bfseries ¿Cuáles son los conceptos clave?}\\[0.16cm]
    \textbf{Redes sociales:} plataformas digitales usadas para comunicación, entretenimiento, información o apoyo académico.\\[0.12cm]
    \textbf{Rendimiento académico:} resultado académico del estudiante, medido mediante el promedio reportado.\\[0.12cm]
    \textbf{Distracción académica:} grado en que el uso de redes puede interrumpir el estudio, tareas o clases.\\[0.12cm]
    \textbf{Hábitos académicos:} horas de estudio y asistencia, variables que ayudan a explicar el desempeño.
};

% 5 Principios
\node[
    anchor=north west,
    text width=8.15cm,
    align=left,
    font=\fontsize{8.4}{10.1}\selectfont
] at (0.7,9.5) {
    {\fontsize{22}{22}\selectfont\bfseries\textcolor{azulUVE}{5}} {\large\bfseries Principios}\\
    {\color{naranjaUVE}\bfseries ¿Cómo sucede el fenómeno?}\\[0.16cm]
    \textbf{Asociación, no causalidad:} los datos permiten observar relaciones estadísticas, pero no demostrar causa directa.\\[0.12cm]
    \textbf{Multifactorialidad:} el rendimiento académico depende de varias variables al mismo tiempo, no solo de las redes sociales.\\[0.12cm]
    \textbf{Uso contextual:} el efecto del uso de redes puede depender de si interfiere o no con las actividades académicas.
};

% 7 Teorias
\node[
    anchor=north west,
    text width=6.45cm,
    align=left,
    font=\fontsize{8.4}{10.1}\selectfont
] at (0.7,5.5) {
    {\fontsize{22}{22}\selectfont\bfseries\textcolor{azulUVE}{7}} {\large\bfseries Teorías}\\
    {\color{naranjaUVE}\bfseries ¿Por qué sucede?}\\[0.16cm]
    \textbf{Desplazamiento del tiempo:} más horas en redes pueden reducir tiempo de estudio o descanso.\\[0.12cm]
    \textbf{Atención limitada:} las interrupciones digitales pueden afectar la concentración.\\[0.12cm]
    \textbf{Uso académico de redes:} algunas plataformas también pueden apoyar el aprendizaje, por lo que el efecto no siempre es negativo.
};

%------------------ LADO DERECHO ------------------

% 4 Procedimiento
\node[
    anchor=north west,
    text width=7.15cm,
    align=left,
    font=\fontsize{8.2}{9.8}\selectfont
] at (19.6,15) {
    {\fontsize{22}{22}\selectfont\bfseries\textcolor{azulUVE}{4}} {\large\bfseries Procedimiento}\\
    {\color{naranjaUVE}\bfseries ¿Qué hicimos para observar el fenómeno?}\\[0.16cm]
    1. Cargar automáticamente la base limpia del proyecto.\\[0.05cm]
    2. Preparar variables: promedio, horas en redes, distracción, estudio y asistencia.\\[0.05cm]
    3. Realizar análisis descriptivo y gráficos exploratorios.\\[0.05cm]
    4. Aplicar prueba t, ANOVA, Kruskal-Wallis y correlaciones.\\[0.05cm]
    5. Ajustar un modelo con horas en redes, distracción, estudio y asistencia.\\[0.05cm]
    6. Revisar diagnósticos del modelo e interpretar resultados.
};

% 6 Registro
\node[
    anchor=north west,
    text width=7.15cm,
    align=left,
    font=\fontsize{8.2}{9.8}\selectfont
] at (19.6,9.45) {
    {\fontsize{22}{22}\selectfont\bfseries\textcolor{azulUVE}{6}} {\large\bfseries Registro y transformación de datos}\\
    {\color{naranjaUVE}\bfseries ¿Qué medimos directamente?}\\[0.16cm]
    \textbf{Fuente:} Student Social Media, Academic Performance Dataset (Kaggle).\\[0.07cm]
    \textbf{Unidad de análisis:} estudiantes universitarios.\\[0.07cm]
    \textbf{Variable respuesta:} promedio académico.\\[0.07cm]
    \textbf{Variables explicativas:} horas en redes, distracción académica, horas de estudio y asistencia.\\[0.07cm]
    \textbf{Transformaciones:} limpieza, recodificación, creación de categorías de uso, tablas, gráficos, pruebas estadísticas y modelo.
};

% 8 Conclusiones
\node[
    anchor=north west,
    text width=7.15cm,
    align=left,
    font=\fontsize{8.15}{9.7}\selectfont
] at (19.6,3.95) {
    {\fontsize{22}{22}\selectfont\bfseries\textcolor{azulUVE}{8}} {\large\bfseries Conclusiones preliminares}\\
    {\color{naranjaUVE}\bfseries ¿Qué podemos afirmar?}\\[0.16cm]
    1. Sí hay evidencia de asociación entre uso de redes sociales y rendimiento académico.\\[0.05cm]
    2. La relación es moderada y no debe interpretarse como causalidad.\\[0.05cm]
    3. Las horas en redes y la distracción se asocian negativamente con el promedio.\\[0.05cm]
    4. Las horas de estudio y la asistencia se asocian positivamente con el promedio.\\[0.05cm]
    5. El rendimiento académico depende de varios factores al mismo tiempo.
};

%------------------ REFERENCIAS ------------------
\node[
    anchor=west,
    text width=23cm,
    align=left,
    font=\footnotesize
] at (0.7,0.00001) {
\textbf{Referencias:}
Kaplan y Haenlein (2010); Chadwick (1979); APA (s.f.); León Luyo et al. (2023); Santos (2010); Aljemely (2020); Ramírez Jaime et al. (2023); Rodelo Moreno y Lizárraga Reyes (2018); Miranda-Rosero et al. (2023).
};

\end{tikzpicture}%
}

\end{document}

